\section{LLM-Based PHP Version Migration Framework}
\label{sec:migration_framework}

We implement a controlled framework for LLM-assisted PHP migration.
It standardizes prompting, context handling, and reconstruction.
The goal is to reduce orchestration variance.
Differences should reflect model behavior, not pipeline artifacts.
Rector is used only for evaluation as the pinned contract.\cite{rector}
The migration procedure is model- and provider-agnostic.

\paragraph{Design principles}
We enforce three constraints.
First, systems share the same degrees of freedom.
We fix prompts, chunking policy, and decoding settings where supported.
Second, extraction and reconstruction are deterministic from recorded responses.
Third, we surface failures.
Timeouts and extraction failures are reported, not repaired.
This isolates model behavior and supports fair comparison.

\subsection{Configuration and provider interface}
\label{sec:framework_config}

A centralized configuration defines model identifiers, request limits, and decoding settings where supported.
We log prompts, responses, and run metadata for audit and recomputation.

A unified provider interface normalizes authentication, request formatting, and error handling.
This keeps the migration policy consistent across systems.

\subsection{Migration policy and context management}
\label{sec:framework_chunking}

We treat context length as a fixed constraint.\cite{liu2024lost,du2025context}
Files up to 500 lines use one request.
Longer files are split to fit a bounded budget.
We use line count as a provider-agnostic proxy for request size.
Chunk boundaries depend only on the input file and pinned parameters.

We use a fixed PHP-aware boundary heuristic to reduce mid-function splits (Appendix~\ref{app:benchmark_algorithms}).
We do not use learned or adaptive chunking.

We disable overlap, retrieval augmentation, multi-pass refinement, and cross-chunk repair.
Each model runs a single-pass policy under the same constraints.

Prompts require begin and end markers.
Whole-file outputs are used directly.
Chunk outputs are concatenated in order.

\subsection{Output extraction and reconstruction}
\label{sec:framework_reconstruction}

Responses must place migrated code between markers.
The parser checks marker presence and extracts only the enclosed code.
For chunked files, we reconstruct outputs by ordered concatenation.

If extraction fails for a chunk, we insert a fixed placeholder comment that reports the failure.
We do not attempt heuristic repair.
Failures propagate to syntax checks and contract scoring.
This prevents inflated results.

We store all artifacts with raw responses and execution metadata.
This supports inspection and reproducibility.

\subsection{Experimental orchestration}
\label{sec:framework_orchestration}

The orchestration layer applies the fixed policy, records extraction outcomes, and prepares outputs for contract evaluation.
This yields consistent experimental conditions and fully traceable runs.

